% GENERATED FILE. Edit canonical lessons, then run npm run generate:books.
% canonical-source-hash: fnv1a64:074e18090cb549a9
% canonical-lessons: GU-C272-himmat, GU-C272-vacan, GU-C272-notarvu, GU-C272-satavvu, GU-C272-asha-rakhvi, GU-R272-review-267-272, GU-R272-review-final-272

\chapter{Promises, hope and courage}
\label{ch:gu-promises-hope-and-courage}
\hlchaptermodality{272}

\begin{chapteropening}
\textbf{By the end of this chapter:} \emph{I can talk about promises, invitations and hope, and say what I dare not do.}

Complete the last lesson of chapter 272: I can talk about promises, invitations and hope, and say what I dare not do.
\end{chapteropening}

\section[\texorpdfstring{\gu{હિંમત} (himmat)}{હિંમત (himmat)}]{\gu{હિંમત} (himmat) — courage}

\label{lesson:GU-C272-himmat}



\begin{quote}
Before the new one: say the Gujarati for foolish, then the Gujarati for stubborn.
\end{quote}

\subsection*{You'll want to know: \gu{હિંમત}}
\textbf{\gu{હિંમત}} — \emph{himmat} — \textquotedblleft{}courage\textquotedblright{}.

\begin{quote}
\textbf{\gu{મારામાં} \gu{બોલવાની} \gu{હિંમત} \gu{નથી}.} — \emph{mārāmā̃ bolvānī himmat nathī} — \textquotedblleft{}I do not have the courage to speak.\textquotedblright{}
\end{quote}

\begin{cousinweb}
\textbf{\gu{હિંમત}} (\emph{himmat}) is a loan from Arabic \emph{himma}, \textquotedblleft{}resolve, aspiration\textquotedblright{}, which reached Gujarati through Persian. Its Arabic root is \emph{h-m-m}, \textquotedblleft{}to intend\textquotedblright{} (Lane).
\end{cousinweb}

\begin{grammarlens}[title={What you've built}]
One more word for everyday life at A2.
\end{grammarlens}

\subsection*{Your turn}
\begin{itemize}
\raggedright
  \item \emph{Say it:} \emph{himmat}
  \item \emph{Say it:} \emph{himmat}, once more
  \item \emph{Say it:} \emph{mārāmā̃ bolvānī himmat nathī} — I do not have the courage to speak.
  \item \emph{Recall:} say the Gujarati for foolish, then the Gujarati for stubborn, then say \emph{himmat} again
\end{itemize}

\begin{tcolorbox}[breakable,colback=teal!4,colframe=teal!35!black,title={Before you move on}]
What does \gu{હિંમત} mean? (\textquotedblleft{}Courage\textquotedblright{}.) Say it once more.
\end{tcolorbox}
\section[\texorpdfstring{\gu{વચન} (vacan)}{વચન (vacan)}]{\gu{વચન} (vacan) — promise, word}

\label{lesson:GU-C272-vacan}



\begin{quote}
Before the new one: say the Gujarati for stubborn, then the Gujarati for courage.
\end{quote}

\subsection*{You'll want to know: \gu{વચન}}
\textbf{\gu{વચન}} — \emph{vacan} — \textquotedblleft{}promise, word\textquotedblright{}.

\begin{quote}
\textbf{\gu{શું} \gu{તમે} \gu{વચન} \gu{આપો} \gu{છો}?} — \emph{śũ tame vacan āpo cho?} — \textquotedblleft{}Do you promise?\textquotedblright{}
\end{quote}

A neuter noun.

\begin{cousinweb}
\textbf{\gu{વચન}} (\emph{vacan}) is Sanskrit \emph{vacaná}, \textquotedblleft{}speaking, statement\textquotedblright{}, taken over as a tatsama. CDIAL also lists the inherited Gujarati \emph{veṇ}, \textquotedblleft{}word, promise\textquotedblright{}, from the same word (CDIAL 11199).
\end{cousinweb}

\begin{grammarlens}[title={What you've built}]
One more word for everyday life at A2.
\end{grammarlens}

\subsection*{Your turn}
\begin{itemize}
\raggedright
  \item \emph{Say it:} \emph{vacan}
  \item \emph{Say it:} \emph{vacan}, once more
  \item \emph{Say it:} \emph{śũ tame vacan āpo cho?} — Do you promise?
  \item \emph{Recall:} say the Gujarati for stubborn, then the Gujarati for courage, then say \emph{vacan} again
\end{itemize}

\begin{tcolorbox}[breakable,colback=teal!4,colframe=teal!35!black,title={Before you move on}]
What does \gu{વચન} mean? (\textquotedblleft{}Promise, word\textquotedblright{}.) Say it once more.
\end{tcolorbox}
\section[\texorpdfstring{\gu{નોતરવું} (notarvũ)}{નોતરવું (notarvũ)}]{\gu{નોતરવું} (notarvũ) — to invite}

\label{lesson:GU-C272-notarvu}



\begin{quote}
Before the new one: say the Gujarati for courage, then the Gujarati for promise.
\end{quote}

\subsection*{You'll want to know: \gu{નોતરવું}}
\textbf{\gu{નોતરવું}} — \emph{notarvũ} — \textquotedblleft{}to invite\textquotedblright{}.

\begin{quote}
\textbf{\gu{શું} \gu{તમે} \gu{પડોશીને} \gu{લગ્નમાં} \gu{નોતર્યા}?} — \emph{śũ tame paḍośīne lagnamā̃ notaryā?} — \textquotedblleft{}Did you invite the neighbours to the wedding?\textquotedblright{}
\end{quote}

\begin{cousinweb}
\textbf{\gu{નોતરવું}} (\emph{notarvũ}) is a tadbhava, a word inherited from Sanskrit and reshaped by sound change. It comes from \emph{nimantrayatē}, \textquotedblleft{}invites\textquotedblright{} (CDIAL 7235).
\end{cousinweb}

\begin{grammarlens}[title={What you've built}]
One more word for everyday life at A2.
\end{grammarlens}

\subsection*{Your turn}
\begin{itemize}
\raggedright
  \item \emph{Say it:} \emph{notarvũ}
  \item \emph{Say it:} \emph{notarvũ}, once more
  \item \emph{Say it:} \emph{śũ tame paḍośīne lagnamā̃ notaryā?} — Did you invite the neighbours to the wedding?
  \item \emph{Recall:} say the Gujarati for courage, then the Gujarati for promise, then say \emph{notarvũ} again
\end{itemize}

\begin{tcolorbox}[breakable,colback=teal!4,colframe=teal!35!black,title={Before you move on}]
What does \gu{નોતરવું} mean? (\textquotedblleft{}To invite\textquotedblright{}.) Say it once more.
\end{tcolorbox}
\section[\texorpdfstring{\gu{સતાવવું} (satāvvũ)}{સતાવવું (satāvvũ)}]{\gu{સતાવવું} (satāvvũ) — to pester, to bother}

\label{lesson:GU-C272-satavvu}



\begin{quote}
Before the new one: say the Gujarati for promise, then the Gujarati for to invite.
\end{quote}

\subsection*{You'll want to know: \gu{સતાવવું}}
\textbf{\gu{સતાવવું}} — \emph{satāvvũ} — \textquotedblleft{}to pester, to bother\textquotedblright{}.

\begin{quote}
\textbf{\gu{મારો} \gu{ભાઈ} \gu{મને} \gu{નથી} \gu{સતાવતો}.} — \emph{māro bhāī mane nathī satāvto} — \textquotedblleft{}My brother does not bother me.\textquotedblright{}
\end{quote}

\begin{cousinweb}
\textbf{\gu{સતાવવું}} (\emph{satāvvũ}) is a tadbhava. It comes from Sanskrit \emph{saṁtāpayati}, \textquotedblleft{}causes to be heated; torments\textquotedblright{}; CDIAL glosses the Gujarati verb \textquotedblleft{}to tease, vex\textquotedblright{} (CDIAL 12886).
\end{cousinweb}

\begin{grammarlens}[title={What you've built}]
One more word for everyday life at A2.
\end{grammarlens}

\subsection*{Your turn}
\begin{itemize}
\raggedright
  \item \emph{Say it:} \emph{satāvvũ}
  \item \emph{Say it:} \emph{satāvvũ}, once more
  \item \emph{Say it:} \emph{māro bhāī mane nathī satāvto} — My brother does not bother me.
  \item \emph{Recall:} say the Gujarati for promise, then the Gujarati for to invite, then say \emph{satāvvũ} again
\end{itemize}

\begin{tcolorbox}[breakable,colback=teal!4,colframe=teal!35!black,title={Before you move on}]
What does \gu{સતાવવું} mean? (\textquotedblleft{}To pester, to bother\textquotedblright{}.) Say it once more.
\end{tcolorbox}
\section[\texorpdfstring{\gu{આશા} \gu{રાખવી} (āśā rākhvī)}{આશા રાખવી (āśā rākhvī)}]{\gu{આશા} \gu{રાખવી} (āśā rākhvī) — to hope}

\label{lesson:GU-C272-asha-rakhvi}



\begin{quote}
Before the new one: say the Gujarati for to invite, then the Gujarati for to pester.
\end{quote}

\subsection*{You'll want to know: \gu{આશા} \gu{રાખવી}}
\textbf{\gu{આશા} \gu{રાખવી}} — \emph{āśā rākhvī} — \textquotedblleft{}to hope\textquotedblright{}.

\begin{quote}
\textbf{\gu{હું} \gu{આશા} \gu{રાખું} \gu{છું} \gu{કે} \gu{તમે} \gu{આવશો}.} — \emph{hũ āśā rākhũ chũ ke tame āvśo} — \textquotedblleft{}I hope that you will come.\textquotedblright{}
\end{quote}

\begin{cousinweb}
\textbf{\gu{આશા} \gu{રાખવી}} (\emph{āśā rākhvī}), to keep a hope, is built from \emph{āśā}, Sanskrit \emph{āś\'{\={a}}}, \textquotedblleft{}hope\textquotedblright{}, taken over as a tatsama (EtymDB; CDIAL 7278 cites it), and the taught verb \emph{rākhvũ}, \textquotedblleft{}to keep\textquotedblright{}.
\end{cousinweb}

\begin{grammarlens}[title={What you've built}]
One more word for everyday life at A2.
\end{grammarlens}

\subsection*{Your turn}
\begin{itemize}
\raggedright
  \item \emph{Say it:} \emph{āśā rākhvī}
  \item \emph{Say it:} \emph{āśā rākhvī}, once more
  \item \emph{Say it:} \emph{hũ āśā rākhũ chũ ke tame āvśo} — I hope that you will come.
  \item \emph{Recall:} say the Gujarati for to invite, then the Gujarati for to pester, then say \emph{āśā rākhvī} again
\end{itemize}

\begin{tcolorbox}[breakable,colback=teal!4,colframe=teal!35!black,title={Before you move on}]
What does \gu{આશા} \gu{રાખવી} mean? (\textquotedblleft{}To hope\textquotedblright{}.) Say it once more.
\end{tcolorbox}
\section[(dialogue)]{People and feelings}

\label{lesson:GU-R272-review-267-272}



\begin{quote}
Say the Gujarati for displeased and disappointed.
\end{quote}

\subsection*{Your turn}
Say these aloud:

\begin{itemize}
\raggedright
  \item displeased, disappointed, to feel restless, to get irritated, to sulk
  \item to quarrel, to scold, to deceive, to persuade, advice
  \item to look after, to bring up, to praise, to remember fondly, respect
  \item temperament, kind, honest, funny, arrogant
  \item polite, generous, selfish, foolish, stubborn
  \item courage, promise, to invite, to pester, to hope
\end{itemize}

\begin{tcolorbox}[breakable,colback=teal!4,colframe=teal!35!black,title={Before you move on}]
Displeased? (\textbf{\gu{નારાજ}}.) Disappointed? (\textbf{\gu{નિરાશ}}.) To feel restless? (\textbf{\gu{અકળાવું}}.) To get irritated? (\textbf{\gu{ચિડાવું}}.) To sulk? (\textbf{\gu{રિસાવું}}.) To quarrel? (\textbf{\gu{લડવું}}.) To scold? (\textbf{\gu{વઢવું}}.) To deceive? (\textbf{\gu{છેતરવું}}.) To persuade? (\textbf{\gu{મનાવવું}}.) Advice? (\textbf{\gu{સલાહ}}.) To look after? (\textbf{\gu{સંભાળવું}}.) To bring up? (\textbf{\gu{ઉછેરવું}}.) To praise? (\textbf{\gu{વખાણવું}}.) To remember fondly? (\textbf{\gu{યાદ} \gu{આવવી}}.) Respect? (\textbf{\gu{આદર}}.) Temperament? (\textbf{\gu{સ્વભાવ}}.) Kind? (\textbf{\gu{દયાળુ}}.) Honest? (\textbf{\gu{પ્રામાણિક}}.) Funny? (\textbf{\gu{રમૂજી}}.) Arrogant? (\textbf{\gu{ઘમંડી}}.) Polite? (\textbf{\gu{નમ્ર}}.) Generous? (\textbf{\gu{ઉદાર}}.) Selfish? (\textbf{\gu{સ્વાર્થી}}.) Foolish? (\textbf{\gu{મૂર્ખ}}.) Stubborn? (\textbf{\gu{જિદ્દી}}.) Courage? (\textbf{\gu{હિંમત}}.) Promise? (\textbf{\gu{વચન}}.) To invite? (\textbf{\gu{નોતરવું}}.) To pester? (\textbf{\gu{સતાવવું}}.) To hope? (\textbf{\gu{આશા} \gu{રાખવી}}.)
\end{tcolorbox}
\section[(dialogue)]{Promises, hope and courage — once more}

\label{lesson:GU-R272-review-final-272}



\begin{quote}
Say the Gujarati for courage and promise.
\end{quote}

\subsection*{Your turn}
\begin{itemize}
\raggedright
  \item \emph{Say it:} courage, promise, to invite, to pester, to hope
\end{itemize}

\begin{tcolorbox}[breakable,colback=teal!4,colframe=teal!35!black,title={Before you move on}]
Courage? (\textbf{\gu{હિંમત}}.) Promise? (\textbf{\gu{વચન}}.) To invite? (\textbf{\gu{નોતરવું}}.) To pester? (\textbf{\gu{સતાવવું}}.) To hope? (\textbf{\gu{આશા} \gu{રાખવી}}.)
\end{tcolorbox}
